% !TeX root = ../本科微分几何笔记.tex
% 由主文档通过 \input 载入；请编译主文档。
\setcounter{chapter}{8} % 保留原书章号。
\chapter{测地曲率、测地线与距离}
\label{chap:dg-geodesics}

\section{Darboux 标架与测地曲率}
设定向曲面内的曲线 $\gamma(s)$ 为单位速，令 $t=\gamma'$、$q=\nu\times t$，则 $(t,q,\nu)$ 为正向标准正交标架，称为\emph{Darboux 标架（Darboux frame）}。加速度分解为
\begin{equation}
 t'=k_gq+k_n\nu,
 \label{eq:dg-darboux-first}
\end{equation}
其中 $k_n=\ip{St}t$ 是法曲率，$k_g=\ip{t'}q$ 称为有符号\emph{测地曲率（geodesic curvature）}。再记 $\tau_g=\ip{St}q$ 为\emph{测地挠率（geodesic torsion）}，有
\[
 q'=-k_gt+\tau_g\nu,\qquad \nu'=-k_nt-\tau_gq.
\]
这些式子由标架内积求导及 $\nu'=-St$ 得到。空间曲率满足 $\kappa^2=k_g^2+k_n^2$；$k_g$ 描述曲线在曲面内部的转弯，$k_n$ 描述曲面造成的外部弯曲。曲线处处沿主方向等价于 $\tau_g=0$。

在上一章的正交标架下写 $t=\cos\theta\,e_1+\sin\theta\,e_2$，则
\begin{equation}
 k_g=\frac{\dd\theta}{\dd s}+\omega(t).
 \label{eq:dg-kg-angle}
\end{equation}
这是测地曲率的切向角表达式。若采用正向正交坐标 $\mathrm I=E\dd u^2+G\dd v^2$，取 $e_1=\mathbf x_u/\sqrt E,e_2=\mathbf x_v/\sqrt G$，由式~\eqref{eq:dg-christoffel} 可得
\[
 \omega=-\frac{E_v}{2\sqrt{EG}}\dd u+
          \frac{G_u}{2\sqrt{EG}}\dd v.
\]
代入便得到常用的 Liouville 测地曲率公式。该式既包括切向量相对标架的转动，又包括标架自身的转动。

\section{测地线方程与初值问题}
\begin{definition}[仿射参数测地线]
曲线 $\gamma(t)$ 若满足 $\nabla_{\dot\gamma}\dot\gamma=0$，称为\emph{测地线（geodesic）}，此时 $t$ 为仿射参数。对非恒定单位速曲线，这等价于 $k_g=0$，即空间加速度完全沿曲面法向。
\end{definition}
度量相容给出 $\dd\norm{\dot\gamma}^2/\dd t=2\ip{\nabla_{\dot\gamma}\dot\gamma}{\dot\gamma}=0$，因此测地线速度恒定。非恒定测地线作任意非线性参数变换后，一般不再满足零协变加速度；仿射变换 $t\mapsto at+b$ 则保持方程。

\begin{theorem}[坐标方程与局部存在唯一性]
坐标曲线 $u(t)=(u^1(t),u^2(t))$ 为仿射参数测地线，当且仅当
\begin{equation}
 \ddot u^k+\sum_{i,j=1}^2\Gamma^k_{ij}(u(t))\dot u^i\dot u^j=0,
 \qquad k=1,2.
 \label{eq:dg-geodesic-equations}
\end{equation}
给定 $p\in\Sigma$ 和 $v\in T_p\Sigma$，存在唯一局部测地线满足 $\gamma(0)=p,\dot\gamma(0)=v$。
\end{theorem}
\begin{proof}
把 $\dot\gamma=\sum_i\dot u^i\mathbf x_i$ 代入式~\eqref{eq:dg-covariant-coordinates} 即得方程。把速度也作为未知量，方程变为光滑的一阶系统；常微分方程初值定理给出局部存在唯一性。不同坐标下得到的解由方程的几何定义和局部唯一性一致拼接。
\end{proof}

\begin{example}[球面与圆柱上的测地线]
半径 $R$ 球面上的单位速测地线满足 $\gamma''=-\gamma/R^2$，因为加速度为法向且 $\ip{\gamma''}\gamma=-1$。因此
\[
 \gamma(s)=p\cos(s/R)+Rv\sin(s/R),\quad
 \norm p=R,\ \norm v=1,\ \ip pv=0.
\]
它位于过球心的固定平面内，故为大圆。圆柱在展开坐标中度量为欧氏度量，测地线满足 $u''=v''=0$；卷回圆柱后得到母线、水平圆或圆柱螺线。其空间曲率未必为零。
\end{example}

\section{能量与长度的第一变分}
对区间 $[a,b]$ 上的曲线定义\emph{能量泛函（energy functional）}
\[
 \mathcal E(\gamma)=\frac12\int_a^b\norm{\dot\gamma}^2\dd t.
\]
与长度不同，能量依赖参数化。Cauchy--Schwarz 不等式给出 $L(\gamma)^2\le2(b-a)\mathcal E(\gamma)$，等号恰在速度恒定时成立。固定端点、固定参数区间的长度极小问题，可在常速参数下转化为能量极小问题。

\begin{theorem}[第一变分公式]
设 $\gamma_\varepsilon(t)$ 是光滑曲线变分，$V=\partial_\varepsilon\gamma_\varepsilon|_0$，则
\begin{equation}
 \left.\frac{\dd\mathcal E(\gamma_\varepsilon)}{\dd\varepsilon}\right|_0
 =\left[\ip V{\dot\gamma}\right]_a^b
 -\int_a^b\ip V{\nabla_{\dot\gamma}\dot\gamma}\dd t.
 \label{eq:dg-energy-variation}
\end{equation}
因此，固定端点的所有光滑变分下能量驻定，当且仅当 $\gamma$ 为仿射参数测地线。若基曲线以弧长参数化且正则，则长度第一变分同样为
\begin{equation}
 \delta L=\left[\ip Vt\right]_a^b-\int_a^b\ip V{\nabla_t t}\dd s.
 \label{eq:dg-length-variation}
\end{equation}
\end{theorem}
\begin{proof}
对能量积分求导，利用度量相容与变分参数的无挠交换关系 $\nabla_{\partial_\varepsilon}\partial_t\gamma=\nabla_{\partial_t}\partial_\varepsilon\gamma$，得 $\delta\mathcal E=\int\ip{\nabla_{\dot\gamma}V}{\dot\gamma}\dd t$。分部积分即得公式。固定端点时 $V(a)=V(b)=0$；任意内部紧支撑切向量场均可通过局部坐标变分实现，因此积分对所有 $V$ 为零恰好迫使协变加速度为零。长度情形对 $\norm{\partial_t\gamma_\varepsilon}$ 求导，并用基曲线速度为 $1$，其余相同。
\end{proof}
驻定是极小的必要条件，不是充分条件。球面上一段大于半个大圆的弧仍是测地线，却有更短的大圆弧连接同样端点。长度驻定还只决定无参数曲线；必须采用常速参数才能写成式~\eqref{eq:dg-geodesic-equations}。

\section{旋转曲面的 Clairaut 关系}
对度量 $\dd s^2+r(s)^2\dd\theta^2$，测地线方程为
\[
 \ddot s-rr'\dot\theta^2=0,\qquad
 \ddot\theta+2\frac{r'}r\dot s\dot\theta=0.
\]
第二式等价于 $r^2\dot\theta=c$。对单位速测地线，若其切向量与母线正方向的有向夹角为 $\alpha$，则 $\sin\alpha=r\dot\theta$，所以
\begin{equation}
 r\sin\alpha=c,\qquad
 \dot s^2+\frac{c^2}{r^2}=1.
 \label{eq:dg-clairaut}
\end{equation}
第一式称为\emph{Clairaut 关系（Clairaut relation）}。它意味着测地线只能进入 $r\ge|c|$ 的区域。母线 $\theta=\text{常数}$ 总是测地线；非恒定的纬线 $s=s_0$ 是测地线的充要条件为 $r'(s_0)=0$。Clairaut 关系是守恒关系，单独使用时仍须结合单位速条件和测地线方程检查解。

\section{指数映射、短程最短性与完备性}
令 $\gamma_v$ 为初值 $(p,v)$ 的测地线，在 $v$ 足够小时它至少定义至 $t=1$。定义\emph{指数映射（exponential map）} $\exp_p(v)=\gamma_v(1)$。解对初值的光滑依赖说明它光滑，且 $\dd(\exp_p)_0$ 为恒等映射：在方向 $w$ 上，$\exp_p(tw)=\gamma_w(t)$，在 $0$ 处求导即得 $w$。逆函数定理遂给出零向量附近的微分同胚，称其像为一个法坐标邻域。

\begin{lemma}[Gauss 引理的二维形式]
在足够小的法坐标邻域去掉中心后，采用切平面极坐标 $(r,\theta)$，则
\begin{equation}
 \mathrm I=\dd r^2+J(r,\theta)^2\dd\theta^2,
 \qquad J(r,\theta)>0.
 \label{eq:dg-geodesic-polar}
\end{equation}
其中 $r$ 方向为从 $p$ 出发的单位速测地线方向，与角度方向正交。
\end{lemma}
\begin{proof}
设 $F(r,\theta)=\exp_p(r e(\theta))$，$e(\theta)$ 为切平面单位圆。径向速度 $F_r$ 的长度恒为 $1$。由无挠性，$\nabla_rF_\theta=\nabla_\theta F_r$，从而
\[
 \partial_r\ip{F_r}{F_\theta}
 =\ip{F_r}{\nabla_\theta F_r}
 =\tfrac12\partial_\theta\norm{F_r}^2=0.
\]
在 $r=0$ 处 $F_\theta=0$，所以该内积为零；局部微分同胚保证 $r>0$ 时 $F_\theta\ne0$。令 $J=\norm{F_\theta}$ 即得结论。
\end{proof}

\begin{theorem}[足够短的测地线实现距离]
每个点 $p$ 都有一个邻域，使从 $p$ 出发的足够短径向测地线在所有连接相同端点的分段光滑曲线中长度最短；最短曲线的像唯一。
\end{theorem}
\begin{proof}
取半径 $b$ 的闭切向圆盘包含于指数映射的微分同胚定义域内。对留在对应开球中的曲线，式~\eqref{eq:dg-geodesic-polar} 给出速度至少为 $|\dot r|$，因此从中心到半径 $a<b$ 的点长度至少为 $a$，径向测地线达到此值。等号要求几乎处处没有角向运动且径向单调，因此路径的像就是径向线段。若竞争曲线离开该开球，在第一次到达边界以前长度已至少为 $b>a$，也不可能更短。曲线经过中心时，可分段使用估计并取端点极限。
\end{proof}
最短性是短程结论；球面上首次到达对跖点时，连接两点的最短大圆弧不再唯一。对一般曲面，指数映射的单射性或非退化性可能在远处失效，需用割点与共轭点理论进一步研究。

\begin{theorem}[Hopf--Rinow 定理：引用]
对连通、无边界、有限维黎曼流形，内蕴距离的完备性等价于每条测地线都可对全部实参数延伸；在这些条件下，任意两点之间存在达到距离的测地线段。
\end{theorem}
该定理作为后续黎曼几何的整体结论引用，不在本书证明。紧致无边界曲面满足完备性，去点平面则不满足。结论保证最短测地线存在，不保证唯一；也不能在未检查完备性时应用。

\section{习题}
\begin{exercise}\label{ex:dg-geodesics-latitude}
半径 $R$ 球面取外法向，余纬角 $0<\theta<\pi$ 的纬线按经度增加方向行进。求其法曲率、测地曲率与空间曲率，并判断何时是测地线。
\end{exercise}
\begin{exercise}\label{ex:dg-geodesics-equator}
证明旋转曲面的纬线成为测地线的条件为 $r'=0$，并求标准环面上的两条这类纬线。
\end{exercise}
\begin{exercise}\label{ex:dg-geodesics-reparameter}
设 $\gamma$ 为非恒定仿射参数测地线，$\phi$ 为区间微分同胚。证明 $\gamma\circ\phi$ 仍为仿射参数测地线的充要条件是 $\phi''=0$。
\end{exercise}
\begin{exercise}\label{ex:dg-geodesics-hyperbolic}
上半平面 $y>0$ 赋予度量 $(\dd x^2+\dd y^2)/y^2$。证明 $K=-1$，验证竖直直线可作测地线，并由测地线方程求出所有非竖直测地线的像。
\end{exercise}

